\chapter{Code Walkthrough}

\section{End-To-End Pipeline}

The runtime path starts when the detector publishes a
\code{vision_msgs/Detection2DArray} on \topic{/detector/armors}. The node does
the following work in one detection callback:

\begin{enumerate}
  \item Confirm camera intrinsics and IMU yaw/pitch have been received.
  \item Convert raw detections into filtered bounding-box candidates.
  \item Run PnP for each candidate.
  \item Transform each successful PnP pose from camera frame to odom frame.
  \item Reject impossible positions by height and distance.
  \item Update the tracker with the remaining armor detections.
  \item Compute an aim result from the tracker state.
  \item Evaluate the fire gate.
  \item Publish the gimbal command, legacy command, HUD pixels, RViz markers,
  and structured debug message.
\end{enumerate}

The key design decision is that the node publishes absolute yaw and pitch in
the microcontroller reference frame. It does not publish relative increments.

\section{\file{src/src/auto_aim_node.cpp}}

\subsection{Responsibility}

\file{auto_aim_node.cpp} is the ROS-facing coordinator. It owns:

\begin{itemize}
  \item all runtime parameters,
  \item subscribers for detections, camera info, and IMU yaw/pitch,
  \item publishers for commands, HUD pixels, markers, and debug data,
  \item the \code{Tracker}, \code{PnPSolver}, \code{FrameTransformer},
  \code{FireGate}, and \code{DebugPublisher} objects,
  \item command smoothing and final sign conversion.
\end{itemize}

It should stay mostly orchestration code. Heavy math belongs in
\file{tracker.cpp}, \file{pnp_solver.cpp}, or \file{frame_transformer.cpp}.

\subsection{Constructor}

The constructor declares parameters and creates all ROS interfaces. Parameter
declaration has two purposes:

\begin{enumerate}
  \item It gives every setting a default value.
  \item It allows a YAML launch file to override that value.
\end{enumerate}

The first large block fills \code{TrackerConfig}. That config is passed into
the tracker constructor:

\begin{lstlisting}[language=C++]
tracker_ = std::make_unique<Tracker>(cfg);
cfg_ = cfg;
\end{lstlisting}

Then the constructor builds a matching \code{FireGate::Config}. Fire-related
parameters are split because the tracker computes planner feasibility while
the fire gate owns the final fire/hold decision.

The target class list is converted into a \code{std::set<std::string>} for
fast lookup. Class \code{"1"} is erased because it represents grey/dead
detections and should not be tracked.

Before starting subscriptions, the constructor validates configuration with
\code{ConfigValidator::validate}. Errors abort startup. Warnings are printed
but allow the node to run. This is deliberate: a wrong sign, impossible bullet
speed, or empty target list is more dangerous than a startup failure.

\subsection{Camera Info Subscription}

The \topic{/camera_info} callback initializes the PnP solver once:

\begin{lstlisting}[language=C++]
std::array<double, 9> K;
std::copy(msg->k.begin(), msg->k.end(), K.begin());
pnp_.init(K, msg->d);
\end{lstlisting}

The matrix $K$ contains focal lengths and principal point. The distortion
vector \code{d} corrects lens distortion. The same callback stores
\code{cam_fx_}, \code{cam_fy_}, \code{cam_cx_}, and \code{cam_cy_} for HUD
pixel projection later.

\subsection{Micro IMU Callback}

\code{microImuCallback} receives \topic{/micro_imu}. The message must contain:

\[
data[0] = yaw,\qquad data[1] = pitch.
\]

Both are in radians and in the microcontroller convention. The node multiplies
them by the configured signs:

\begin{lstlisting}[language=C++]
imu_yaw_   = yaw_sign_ * micro_yaw_rad;
imu_pitch_ = pitch_sign_ * micro_pitch_rad;
\end{lstlisting}

Then it updates the frame transformer. This creates the cached rotation and
translation used by detection callbacks. If no IMU message has arrived, the
detection callback refuses to run because camera points cannot be transformed
into the correct frame.

\subsection{Detection Callback}

\code{detectionCallback} is the main hot path. It starts with two guards:

\begin{itemize}
  \item PnP must be initialized from camera info.
  \item IMU yaw/pitch must be valid.
\end{itemize}

It creates a \code{DebugFrame} at the start and fills it stage by stage. This
is why the debug topic can show where a frame failed.

Frame time \code{dt} is computed from message timestamps and clamped to
$0.01$ to $0.10$ seconds. The clamp prevents one bad timestamp from exploding
the EKF prediction.

The detection adapter filters raw detections:

\begin{lstlisting}[language=C++]
auto candidates = DetectionAdapter::convert(*msg, target_classes_);
\end{lstlisting}

For each candidate, the node runs PnP:

\begin{lstlisting}[language=C++]
auto pnp = pnp_.solve(cand.cx, cand.cy, cand.w, cand.h,
                      cfg_.light_ratio, 10.0, enable_pnp_refine_);
\end{lstlisting}

Failed PnP results are still written into the debug frame when possible. That
helps diagnose whether a detection was dropped because of PnP quality rather
than tracker logic.

Successful PnP results go through the frame transformer:

\begin{lstlisting}[language=C++]
auto td = frame_transformer_.apply(pnp.rvec, pnp.tvec, max_oblique_rad);
\end{lstlisting}

Then the node rejects positions with excessive height or range:

\begin{itemize}
  \item \param{max_armor_z} limits vertical position.
  \item \param{max_armor_distance} limits horizontal range.
\end{itemize}

The remaining detections are converted into \code{ArmorDetection} values and
passed to the tracker:

\begin{lstlisting}[language=C++]
tracker_->update(armors, dt);
\end{lstlisting}

After the tracker update, the node copies EKF state, innovation, Mahalanobis
distance, and target id into the debug frame.

\subsection{Aim Computation And Latency}

The node calls:

\begin{lstlisting}[language=C++]
auto aim = tracker_->computeAim(imu_yaw_, imu_pitch_, override_pred_lead);
\end{lstlisting}

Normally \code{override_pred_lead} is \code{-1.0}, which tells the tracker to
use the legacy \param{time_bias} calculation. If \param{use_measured_latency}
is enabled and initialized, the node passes an EMA of measured pipeline latency
plus \param{gimbal_response_delay_s}.

The callback also projects the selected impact point back into camera-relative
yaw and pitch. The fire gate uses this camera-frame total angle as the
\code{OFF_AXIS} check.

\subsection{Command Publishing}

\code{publishCommand} is where the final command is formed. The order matters:

\begin{enumerate}
  \item Start from \code{aim.abs_yaw} and \code{aim.abs_pitch}.
  \item Subtract bore-sight offsets.
  \item Store pre-smoothing values for debug.
  \item Apply EMA smoothing if tracking.
  \item Compute smoothing lag.
  \item Call \code{FireGate::evaluate}.
  \item Convert yaw and pitch back to the microcontroller convention using
  \param{gimbal.yaw_sign} and \param{gimbal.pitch_sign}.
  \item Publish \code{geometry_msgs::Twist}.
\end{enumerate}

The primary command contract is:

\begin{longtable}{p{0.24\textwidth}p{0.62\textwidth}}
\toprule
Field & Meaning \\
\midrule
\code{angular.z} & absolute yaw target in radians \\
\code{angular.y} & absolute pitch target in radians \\
\code{angular.x} & fire flag, $1.0$ means fire and $0.0$ means hold \\
\code{linear.x} & range to selected armor face in meters \\
\bottomrule
\end{longtable}

The same yaw, pitch, and fire values are also published on the legacy topic
\topic{/cmd_vel_AI}. The legacy topic does not include distance.

\subsection{Markers}

\code{publishMarkers} publishes RViz markers:

\begin{itemize}
  \item enemy center sphere,
  \item four armor cubes,
  \item velocity arrow,
  \item approximate aim point,
  \item selected impact point.
\end{itemize}

When the tracker is not tracking, the function publishes delete actions for
old markers so RViz does not show stale geometry.

\section{\file{src/include/auto_aim/debug_frame.hpp}}

\code{DebugFrame} is an internal plain data struct. The detection callback
fills it during processing. \code{DebugPublisher} copies it into the ROS
message at the end.

This two-step design keeps the hot callback readable. Each stage can write its
own fields without constructing a ROS message directly.

\section{\file{src/src/detection_adapter.cpp}}

\subsection{Responsibility}

\code{DetectionAdapter::convert} turns a
\code{vision_msgs::msg::Detection2DArray} into a vector of
\code{DetectionCandidate} structs.

It performs three checks:

\begin{enumerate}
  \item Ignore detections with no classification result.
  \item Ignore detections whose class id is not in \param{target_classes}.
  \item Ignore boxes smaller than \code{min_pixel_size}.
\end{enumerate}

The adapter does not know about PnP, tracking, or firing. That separation is
useful: detector-message parsing can be tested and changed without touching
the EKF.

\section{\file{src/src/pnp_solver.cpp}}

\subsection{Initialization}

\code{PnPSolver::init} stores camera intrinsics and distortion. It also builds
two 3D armor models:

\begin{itemize}
  \item small armor,
  \item large armor.
\end{itemize}

The object points lie on a flat plane with $x=0$ and dimensions converted from
millimeters to meters.

\subsection{Solving}

\code{PnPSolver::solve} first rejects invalid input:

\begin{itemize}
  \item non-finite center or size,
  \item boxes below minimum size,
  \item invalid \param{light_ratio}.
\end{itemize}

It then builds four synthetic image points from the bounding box. The helper
\code{try_model} runs OpenCV \code{solvePnP} with \code{SOLVEPNP_IPPE}, which
is suited for planar targets. If \param{enable_pnp_refine} is true, it runs a
Levenberg-Marquardt refinement step.

Both small and large armor models are tested. The lower reprojection error
wins. The output \code{PnPResult} contains:

\begin{itemize}
  \item \code{ok}: true only if solve succeeded and error is below threshold,
  \item \code{solver_ok}: true if PnP itself produced a finite result,
  \item \code{is_large}: which model won,
  \item \code{rvec} and \code{tvec}: OpenCV pose output,
  \item \code{reproj_err} and \code{reproj_err_norm}: quality metrics,
  \item \code{image_points}: synthetic corners used for debugging.
\end{itemize}

\section{\file{src/src/frame_transformer.cpp}}

\subsection{updateFromImu}

\code{updateFromImu} converts current yaw and pitch into a camera-to-odom
rotation.

\begin{lstlisting}[language=C++]
q_gimbal.setRPY(0.0, imu_pitch, imu_yaw);
q_conv.setRPY(-M_PI / 2.0, 0.0, -M_PI / 2.0);
rotation_ = q_gimbal * q_conv;
translation_ = tf2::Vector3(0, 0, gimbal_height_);
\end{lstlisting}

\code{q_conv} is the fixed rotation from gimbal body axes to camera optical
axes. \code{q_gimbal} changes with the IMU. The translation lifts the camera
to gimbal height.

\subsection{apply}

\code{apply} receives OpenCV \code{rvec} and \code{tvec}. It:

\begin{enumerate}
  \item converts \code{tvec} into a camera-frame point,
  \item converts \code{rvec} into a quaternion,
  \item rotates and translates position into odom,
  \item rotates orientation into odom,
  \item extracts yaw,
  \item corrects or clamps impossible armor yaw,
  \item returns \code{TransformedDetection}.
\end{enumerate}

If orientation extraction fails, \code{valid} remains false and the detection
is ignored.

\section{\file{src/src/tracker.cpp}}

\subsection{Responsibility}

\code{Tracker} is the largest math module. It owns:

\begin{itemize}
  \item EKF state \code{x_},
  \item covariance \code{P_},
  \item target state machine,
  \item target association,
  \item armor face jump handling,
  \item radius and alternate-radius tracking,
  \item ballistic aim computation,
  \item optional adaptive process noise and anti-gyro timing.
\end{itemize}

\subsection{Constructor}

The constructor initializes a 9-element zero state and a diagonal initial
covariance \code{P0_}. Larger covariance means less certainty. For example,
yaw rate starts with relatively large uncertainty because it is hard to know
from one detection.

\subsection{ekfPredict}

\code{ekfPredict} runs before measurement association unless the tracker is
\code{LOST}. It applies damped constant velocity to position, height, and yaw.

When state is \code{TEMP_LOST}, it uses \param{alpha_coast} to decay velocity
faster because no measurement is currently confirming the motion.

If \param{enable_adaptive_q} is true, process noise is changed dynamically:

\begin{itemize}
  \item stationary targets get reduced process noise to suppress jitter,
  \item high yaw-rate targets get increased yaw process noise for response.
\end{itemize}

After prediction, covariance is symmetrized and bounded. Bounding prevents
coasting uncertainty from growing without limit.

\subsection{ekfUpdate}

\code{ekfUpdate} applies a measurement
\[
\vect{z} = [x_a, y_a, z_a, yaw]^T.
\]

It builds the Jacobian \code{H}, the range- and obliquity-dependent
measurement noise \code{R}, then computes innovation, Kalman gain, state
update, and Joseph-form covariance update.

Two debug values are stored:

\begin{itemize}
  \item \code{last_innovation_norm_},
  \item \code{last_mahalanobis_}.
\end{itemize}

These are published so operators can tell whether the tracker agrees with
incoming detections.

\subsection{ekfMahalanobis}

This function computes the association distance for a candidate detection. A
candidate must have distance below \param{maha_threshold} to be associated
with the active target.

If the math solve fails, it returns a very large number so the detection is
rejected.

\subsection{initFromDetection}

When the tracker is \code{LOST}, the closest valid detection initializes the
state. The detected armor is converted into an estimated enemy center by
adding radius in the face-yaw direction:

\[
x_c = x_a + r\cos\psi,\qquad y_c = y_a + r\sin\psi.
\]

Velocities start at zero. The target id becomes the detection's class id.

\subsection{shouldSwitch}

Target switching is conservative. The tracker does not switch while
\code{DETECTING} or \code{TRACKING}. In \code{TEMP_LOST}, it can switch after a
cooldown if a new target is clearly closer according to
\param{switch_range_ratio}.

This avoids bouncing between two similar targets.

\subsection{handleArmorJump}

When the target rotates and another face becomes visible, measured yaw can
jump. \code{handleArmorJump} snaps yaw, handles possible spin reversal, updates
alternate face radius/height information, and either resets or inflates
covariance depending on whether the inferred armor position is plausible.

Then it runs a normal EKF update using the jump measurement.

\subsection{update}

\code{update} is the main per-frame tracker entry point. It has these stages:

\begin{enumerate}
  \item Decrement target-switch cooldown.
  \item Possibly switch targets if state allows it.
  \item Initialize from closest detection if state is \code{LOST}.
  \item Predict EKF state.
  \item Associate detections using class id and Mahalanobis gate.
  \item Run normal update, face-jump handling, or miss handling.
  \item Clamp radius and yaw rate.
  \item Advance the state machine.
\end{enumerate}

The state machine logic at the end is important:

\begin{itemize}
  \item \code{DETECTING} needs \param{confirm_frames} consecutive matches.
  \item \code{TRACKING} becomes \code{TEMP_LOST} after one miss.
  \item \code{TEMP_LOST} becomes \code{LOST} after \param{lost_timeout}.
\end{itemize}

\subsection{solveBallistic}

\code{solveBallistic} returns pitch and flight time for a target at horizontal
distance \code{gd} and vertical difference \code{dz}. It iterates five times
to compensate gravity. It rejects too-close, too-slow-horizontal, or extreme
pitch solutions.

\subsection{computeAim}

\code{computeAim} turns the EKF state into a gimbal command. It returns an
empty \code{AimResult} unless the tracker is \code{TRACKING} or
\code{TEMP_LOST}.

The function computes the barrel position from current yaw and configured
barrel offset. It then evaluates four future faces over two passes. Each face
candidate stores:

\begin{itemize}
  \item predicted 3D face position,
  \item range and bearing from barrel,
  \item ballistic pitch and flight time,
  \item fire margin,
  \item visibility score,
  \item anti-gyro residual when applicable.
\end{itemize}

The best face is the one with highest margin, with visibility as a tie-break.
In anti-gyro mode, the best viable residual can override the margin pick.

The final \code{AimResult} contains:

\begin{itemize}
  \item absolute yaw and pitch,
  \item relative yaw and pitch for debug/fire only,
  \item distance,
  \item selected target point,
  \item tracking and fire booleans,
  \item selected face index,
  \item flight time and prediction time,
  \item anti-gyro residual.
\end{itemize}

\section{\file{src/src/fire_gate.cpp}}

\subsection{Responsibility}

\code{FireGate} is the single source of truth for fire/hold. It receives
\code{FireGate::Inputs} from the command publisher and returns a
\code{Decision} with:

\begin{itemize}
  \item \code{fire}: true or false,
  \item \code{blocker}: enum reason,
  \item \code{reason}: human-readable debug string.
\end{itemize}

\subsection{Decision Order}

The order is deliberate:

\begin{enumerate}
  \item Not tracking blocks immediately.
  \item Invalid target blocks immediately.
  \item Invalid ballistic solution blocks immediately.
  \item Out-of-range blocks immediately.
  \item Optional stale measurement blocks.
  \item Optional anti-gyro timing blocks.
  \item Margin, off-axis, and smoothing checks decide normal fire.
  \item Hysteresis may keep fire true across small angular drift.
  \item The most informative remaining blocker is returned.
\end{enumerate}

Hard failures reset hysteresis. Hysteresis only helps small geometry-like
drifts after fire was already allowed.

\section{\file{src/src/debug_publisher.cpp}}

\code{DebugPublisher} creates the \topic{/auto_aim/debug} publisher. Its
\code{publish} function copies every \code{DebugFrame} field into the generated
ROS message type.

It also maintains a fire blocker histogram. Every
\code{fire_log_period_s} seconds it logs counts and percentages for recent
blocker reasons. This is useful when fire is not allowed and individual
messages are too noisy to inspect by eye.

\section{\file{src/src/config_validator.cpp}}

\code{ConfigValidator} checks parameter ranges before the node fully starts.
It separates warnings from errors:

\begin{itemize}
  \item Errors mean the node should not run. Examples: empty target classes,
  impossible sign values, or max fire distance less than min fire distance.
  \item Warnings mean the value is unusual but not necessarily impossible.
  Examples: nonzero bore-sight offsets or suspicious physical dimensions.
\end{itemize}

Use the validator as the first place to add checks when adding new parameters.
Bad parameter values should fail loudly at startup, not silently mis-aim on
the field.

\section{Headers In \file{src/include/auto_aim}}

The headers expose the package's internal module boundaries:

\begin{longtable}{p{0.34\textwidth}p{0.56\textwidth}}
\toprule
Header & What to read there \\
\midrule
\file{tracker.hpp} & Main data structs, tracker config, public tracker API,
state enum, and private helper declarations. \\
\file{pnp_solver.hpp} & PnP result fields, armor dimensions, solve signature. \\
\file{frame_transformer.hpp} & Frame definitions and transformed detection
output. \\
\file{fire_gate.hpp} & Fire blocker enum, input fields, config fields, and
decision output. \\
\file{detection_adapter.hpp} & Detection candidate representation and convert
function. \\
\file{debug_frame.hpp} & Internal debug fields populated across the callback. \\
\file{debug_publisher.hpp} & Debug publisher interface. \\
\file{config_validator.hpp} & Validation result and validator interface. \\
\bottomrule
\end{longtable}

\section{Messages}

\subsection{\file{src/msg/AutoAimDebug.msg}}

This is the most important message for debugging. It groups fields by pipeline
stage:

\begin{itemize}
  \item detection,
  \item PnP,
  \item frame transform,
  \item EKF,
  \item aim planner,
  \item command publisher,
  \item fire gate,
  \item latency.
\end{itemize}

When a bug appears, inspect these groups in order. For example, if
\code{pnp_tvec_z} is unstable on a static target, do not tune the EKF first.

\subsection{\file{src/msg/GimbalCmd.msg}}

This message is not used at runtime. The actual command topic is
\topic{/tracker/cmd_gimbal} with \code{geometry_msgs/Twist}. Keep
\file{GimbalCmd.msg} only as a placeholder unless the team intentionally
migrates the firmware contract.

\section{Configuration Files}

\file{params_realsense_16.yaml} and \file{params_zed_64.yaml} provide
robot-specific runtime values. The two files are similar, but physical values
such as \param{gimbal_height}, \param{barrel_offset_x}, and
\param{barrel_offset_z} should be measured per robot.

Important categories:

\begin{itemize}
  \item detection/tracking: target classes, light ratio, max armor range,
  state machine settings;
  \item EKF: process noise, measurement noise, damping;
  \item physical: bullet speed, gravity, gimbal height, barrel offset;
  \item fire gate: angular window, range limits;
  \item timing: prediction bias, reference frequency;
  \item command shaping: smoothing alpha and sign convention.
\end{itemize}

If the code has a default value and the YAML has a different value, the YAML
wins at launch.

\section{Launch Files}

\file{auto_aim_realsense_16.launch.py} and \file{auto_aim_zed_64.launch.py}
load the matching YAML file and start \code{auto_aim_node}. They assume
upstream nodes provide:

\begin{itemize}
  \item \topic{/detector/armors},
  \item \topic{/camera_info},
  \item \topic{/micro_imu}.
\end{itemize}

\file{auto_aim_launch.py} declares parameters inline. It is useful as an older
single-file example, but the robot-specific YAML launch files are clearer for
real deployments.

\section{Build Files}

\subsection{\file{src/CMakeLists.txt}}

CMake does the following:

\begin{enumerate}
  \item sets C++17,
  \item finds ROS 2, OpenCV, Eigen, tf2, and message-generation dependencies,
  \item generates message code from \file{msg/GimbalCmd.msg} and
  \file{msg/AutoAimDebug.msg},
  \item builds \code{libauto_aim_component.so} from all C++ source files,
  \item links generated message typesupport,
  \item registers \code{auto_aim::AutoAimNode} as a composable node and
  standalone executable,
  \item installs headers, launch files, config files, and docs.
\end{enumerate}

If a new C++ source file is added, it must be listed in
\code{add_library(...)} or it will not be compiled.

\subsection{\file{src/package.xml}}

\file{package.xml} is ROS package metadata. It names dependencies so ROS tools
know what must be present to build and run the package.

\section{Debugging Recipes By Symptom}

\subsection{Node Starts But Publishes No Commands}

Check:

\begin{enumerate}
  \item Is \topic{/camera_info} publishing?
  \item Is \topic{/micro_imu} publishing at least two floats?
  \item Is \topic{/detector/armors} publishing?
  \item Does \topic{/auto_aim/debug} show \code{detection_present=true}?
  \item Is tracker state stuck in \code{LOST} or \code{DETECTING}?
\end{enumerate}

\subsection{PnP Fails Often}

Check:

\begin{enumerate}
  \item Bounding boxes are not too small.
  \item \param{light_ratio} is reasonable.
  \item Camera intrinsics are correct.
  \item \code{pnp_reproj_err_norm} is not large on clean frames.
\end{enumerate}

\subsection{Target Position Jitters While Robot Is Still}

Check in order:

\begin{enumerate}
  \item bbox center and size jitter,
  \item PnP \code{tvec},
  \item \topic{/micro_imu} yaw and pitch noise,
  \item odom position,
  \item EKF innovation and covariance behavior.
\end{enumerate}

\subsection{Pitch Is Consistently Wrong}

Check:

\begin{enumerate}
  \item measured \param{bullet_speed},
  \item \param{gimbal_height},
  \item \param{barrel_offset_z},
  \item pitch sign,
  \item whether \param{pitch_offset_deg} was used to hide a calibration issue.
\end{enumerate}

\subsection{Fire Is Never Allowed}

Check \code{fire_blocker} and \code{fire_blocker_reason}. Common causes:

\begin{itemize}
  \item \code{NOT_TRACKING}: tracker has not reached \code{TRACKING};
  \item \code{OUT_OF_RANGE}: distance outside configured range;
  \item \code{MARGIN_NEGATIVE}: selected face is not aligned;
  \item \code{OFF_AXIS}: target is outside camera angular window;
  \item \code{SMOOTHING_LAG}: command smoothing has not caught up.
\end{itemize}

\section{Good First Code Tasks}

\begin{enumerate}
  \item Add a new field to \code{DebugFrame} and \code{AutoAimDebug.msg}, then
  publish it through \code{DebugPublisher}. This teaches the full debug-message
  path.
  \item Add a validator warning for a new parameter. This teaches safe startup
  checks.
  \item Add a console log for one state transition in \code{Tracker::update}.
  This teaches tracker flow without changing math.
  \item Add a small offline unit test for \code{FireGate::evaluate}. This is a
  contained logic module with no ROS subscriptions.
\end{enumerate}

\section{Rules For Future Maintainers}

\begin{itemize}
  \item Keep units visible in comments and parameter names.
  \item Keep frame conversions inside \code{FrameTransformer} or very close to
  where data enters/leaves a frame.
  \item Keep fire decision logic centralized in \code{FireGate}.
  \item Do not change the \topic{/tracker/cmd_gimbal} contract without
  firmware coordination.
  \item Prefer adding debug fields over guessing from terminal logs.
  \item When fixing a bug, identify the earliest pipeline stage where the data
  becomes wrong.
\end{itemize}
